% PAPER_TODO 2.2: evaluation with MEASURED numbers only.  Every number has a source comment (file in the repository);
% docs/paper/NUMBERS_MAP.md maps each projected number of the draft to these, or marks it as dropped.
% Framing-neutral: usable for frame (A) theory+method or (B) oracle study (PAPER_TODO 0.2); reorder subsections accordingly.
\section{Evaluation}
\label{sec:evaluation}

All studies were pre-registered: protocol, configuration and tools were hashed before any outcome was seen, seeds were drawn
from the public drand beacon after the freeze, and amendments are recorded and labelled.  Results that were analysed after
seeing data are marked as such.  Every number below is produced by a workflow in the artifact.

\subsection{Where simple gates fail: the phase diagram}
\label{sec:eval-phase}
% source: paper_snippets/P0_AND_PHASE_DIAGRAM.md, tools/gate_phase_diagram.py; study_h/results/RESULTS_SUMMARY_study-h-v1.md
For a patch that always helps, Table~\ref{tab:phase} gives the probability that each gate admits it as a function of the true
harm rate $h$ ($\eta=5\%$, $\alpha=0.1$; exact for the smoke, zero-tolerance and fixed-sample gates, Monte Carlo for the
sequential rules).  A zero-tolerance gate with $300$ draws keeps fewer than half of the acceptable patches once $h>0.23\%$; a
$20$-draw smoke test still admits more than $10\%$ of inadmissible patches until $h\approx11\%$.  Outside the band
$[3/n,\,1-0.1^{1/20}]\approx[0.75\%,10.9\%]$ ($n=400$) a simple gate suffices.  At the same number of draws a fixed-sample test is
as safe as the sequential certificate and \emph{more} powerful ($0.92$ versus $0.65$ at $h=2\%$, $300$ draws).  What the
sequential rule buys is validity at data-dependent stopping and early stopping (110 clean draws in a 16-cell family), not power.
On synthetic targets with rare harm (Study H, $50$ repetitions per cell, $300$ draws), the certificate admitted $0/100$ cells with
$\rho\ge\eta$, while a $5$-draw smoke test admitted $38/50$ and $27/50$ at $\rho=5\%$ and $10\%$, and replay plus fingerprints
admitted every cell.

\begin{table}[t]
\centering\footnotesize
\caption{Probability that a gate admits an always-helpful patch, by true harm rate $h$ ($h\ge\eta$: admitting is an error).}
\label{tab:phase}
\begin{tabular}{@{}lccccc@{}}
\toprule
$h$ & smoke 20 & zero-tol.\ 300 & fixed 300 & certificate, 300 & mixture, 1000 \\
\midrule
0.3\%  & 0.94 & 0.41 & 1.00 & 1.00 & 1.00 \\
1\%    & 0.82 & 0.05 & 1.00 & 0.95 & 1.00 \\
2\%    & 0.67 & 0.00 & 0.92 & 0.65 & 1.00 \\
3\%    & 0.54 & 0.00 & 0.59 & 0.36 & 0.78 \\
5\% (error)  & 0.36 & 0.00 & 0.06 & 0.10 & 0.10 \\
10\% (error) & 0.12 & 0.00 & 0.00 & 0.01 & 0.00 \\
\bottomrule
\end{tabular}
\end{table}

\subsection{Skill reuse on real targets}
\label{sec:eval-skills}
% sources: study_d/results/RESULTS_SUMMARY_study-d-v1.md; study_e/results/RESULTS_SUMMARY_study-e-v1.md;
%          study_g/results/RESULTS_SUMMARY_study-g-v1.md; harness_checks/STATELESSNESS_SUMMARY.md
Two skill families, both frozen AST rewrites (excluding \texttt{bool} from \texttt{int} checks, learned from one PySStuBs fix, and
\texttt{== None}$\to$\texttt{is None}),
were applied to $14$ real targets in three pre-registered studies ($28$ cells).
\emph{Study D} ($8$ targets, $16$ cells): ground truth from $300$ draws plus the target's full test suite gave $3$ beneficial,
$9$ unchanged, $4$ out-of-scope and $0$ harmful cells.  B2, replay plus fingerprints (B6-R) and always-retest (B8) each admitted
$12$ cells, the $3$ beneficial ones and all $9$ unchanged ones; the certificate (benefit e-process with e-Holm over the $16$ cells, family-wise $\alpha=0.1$) admitted $2$ cells, both
beneficial, after $110$ draws each, and missed the third, whose mean benefit ($-0.11$) is smaller than the margin $\theta=0.2$.  It used $3{,}220$ trial runs
against $3{,}600$ for B8.
\emph{Studies E and G}: on the four harmful cells (\texttt{pyyaml}, \texttt{packaging}, \texttt{databricks}: harm on $87$--$100\%$ of
inputs; \texttt{pydal}: $81/300$), B2 and B6-R admitted all four and the certificate rejected all four, within two draws for the first
three.  The targets' own test suites ($1{,}287$, $51{,}524$ and $130$ tests) showed \emph{no} new regression in Study E; in Study G
they did ($21$ new failures).
Re-running $2{,}400$ recorded draws in shuffled order with fresh workers reproduced every paired outcome, so the draws behave as
independent, as Assumption~\ref{ass:trials} requires.
\emph{Reading.} On these targets harm was frequent and deterministic, so a simple gate with target evidence suffices; the
certificate's distinctive behaviour is refusing changes that are not beneficial, not catching rare harm.

\subsection{Partial harm in LLM patches}
\label{sec:eval-patches}
% sources: research_closure/results/study_p/P_RESULT.md; research_closure/results/rq5/RQ5_RESULT.md
Study P drew fresh repository changes from a pre-registered frame of $3{,}562$ candidates; for each usable instance a light
model generated one patch, and $205$ patches were plausible (they pass the hidden tests).  With a function-level oracle ($400$ inputs per patch), $16/205=7.8\%$ (Wilson $4.9$--$12.3\%$) had an
estimated regression rate inside the partial-harm band.  The test-based census admitted all $16$ patches with $\hat\theta>5\%$, a
$20$-draw smoke test admitted $2$, and the zero-tolerance, fixed-sample and sequential gates admitted none.  A pre-registered manual
check (RQ5) then found that $97.5\%$ of the ``regression'' inputs were inputs the developers' fix also rejects; counting only
regressions on inputs the fix accepts left $1$ of $205$ patches as planned, $3$ when deliberate rejections that the patch skips are
added (after seeing data), and at most $6$ when non-reproducible patches are counted.

\subsection{Naive differences versus contract regressions on SWE-bench Verified}
\label{sec:eval-swebench}
% source: research_closure/results/study_q/STUDY_Q_RESULT.md; qa1/QA1_RESULT.json
Study Q tested the RQ5 finding on an independent, widely used benchmark.  For $469$ SWE-bench Verified instances with at least one
public resolved patch, up to three distinct resolved patches were scored in the official Docker images ($245$ instances usable,
$694$ patches).  $43.1\%$ of patches behaved differently from the developers' fix on some input; $11.1\%$ (Wilson $9.0$--$13.6\%$)
had a \emph{contract regression}, an input on which the fix returns a value and the patch does not, or the fix rejects on purpose
and the patch accepts.  The confirmatory hypothesis, that most naive differences are not contract regressions, was \emph{not}
supported: $6/55$ ($10.9\%$, Wilson $5.1$--$21.8\%$; one-sided $p\approx1$).
A manual review added after the result (labelled as such) re-ran the $55$ units with outputs kept: $36$ reproduced exactly, and two
independent raters agreed on $87.7\%$ of $154$ items.  $8/36$ ($22.2\%$, Wilson $11.7$--$38.1\%$) units had no regression confirmed by
both raters; excluding three units whose inputs are model-built stand-in objects gives $11/36$ ($18.0$--$46.9\%$).  The RQ5 conclusion
therefore does not generalise: on SWE-bench Verified most naive differences are real, observable behaviour differences from the fix.
Whether they are \emph{wrong} is a separate question that a reference implementation cannot answer.

\subsection{Who defines harm}
\label{sec:eval-oracle}
% source: research_closure/results/rq3/RQ3_RESULT.md
On $14$ measurable targets of Studies D, E and G, a model-written input generator agreed with the authors' hand-written oracles on
$11/14$ detection decisions (Cohen's $\kappa=0.57$), a signature-based generator on $8/14$ ($\kappa=0.19$).  In one disagreement the
generated inputs were right: on \texttt{cerberus} the skill rewrote two functions other than the intended one; the hand-written
oracle, which checks the intended region, saw nothing, while generated inputs changed behaviour on $24$--$56$ of $200$ inputs.  An
oracle written from the author's intent misses side effects outside that intent, which is the scope problem this paper is about.
